\subsection{应用：正项级数的第二类收敛判别法}

我们经常遇到这样的级数 \((u_{n})\) ：从某项起 \(u_{n} > 0\) ，且 \(u_{n+1}/u_{n}\) 有一个容易确定的渐近展开。对这样的级数，拥有一些（称为第二类判别法的）判别法，使得仅从 \(u_{n+1}/u_{n}\) 的形状就能确定级数是否收敛，是很方便的。下面就是这样一个判别法：

命题 7（“拉阿贝判别法”）. 设 \((u_{n})\) 是从某项起各项 \textgreater{} 0 的级数。若从某个时刻起 \(u_{n+1}/u_{n} \leqslant 1 - \frac{\alpha}{n}\) （对某个 \(\alpha > 1\) ），则级数 \((u_{n})\) 收敛；若从某个时刻起有 \(u_{n+1}/u_{n} \geqslant 1 - \frac{1}{n}\) ，则级数 \((u_{n})\) 的和为无穷。

事实上，若 \(u_{n+1} / u_n \leqslant 1 - \frac{\alpha}{n}\) （其中 \(\alpha > 1\) ）对一切 \(n \geqslant n_0\) 成立，则有 \(u_n \preccurlyeq p_n = \prod_{k=n_0}^{n} \left(1 - \frac{\alpha}{k}\right)\) 。而 \(\log \left(1 - \frac{\alpha}{n}\right) = -\frac{\alpha}{n} - \frac{\alpha^2}{2n^2} + o\left(\frac{1}{n^2}\right)\) ，由此 \(\log p_n =\) \(-\alpha \log n + k + o(1/n) (k \text{ 为常数}), \text{ 且 } p_n \sim e^k \frac{1}{n^\alpha}; \text{ 由于 } \alpha > 1 \text{ 0 阶对数判别法足以完成证明。}\)

反之，若从某个时刻起有 \(u_{n + 1} / u_n \geqslant 1 - \frac{1}{n}\) ，则同样的计算证明 \(u_n \succcurlyeq \frac{1}{n}\) ，命题得证。

利用阶 \textgreater{} 0 的对数判别法，同样可以证明下列第二类判别法：

命题 8. 设 \((u_n)\) 是从某个时刻起各项 \(>0\) 的级数。若从某个时刻起有

\[
\frac {u _ {n + 1}}{u _ {n}} \leqslant 1 - \frac {1}{n} - \frac {1}{n l _ {1} (n)} - \dots - \frac {1}{n l _ {1} (n) l _ {2} (n) \dots l _ {p - 1} (n)} - \frac {\alpha}{n l _ {1} (n) l _ {2} (n) \dots l _ {p} (n)}
\]

（其中 \(\alpha > 1\) ），则级数 \((u_n)\) 收敛；若从某个时刻起有

\[
\frac {u _ {n + 1}}{u _ {n}} \geqslant 1 - \frac {1}{n} - \frac {1}{n l _ {1} (n)} - \dots - \frac {1}{n l _ {1} (n) l _ {2} (n) \dots l _ {p} (n)}
\]

则级数 \((u_{n})\) 的和为无穷。

例. 考察超几何级数，其通项为

\[
u _ {n} = \frac {\alpha (\alpha + 1) \dots (\alpha + n - 1) \beta (\beta + 1) \dots (\beta + n - 1)}{1 . 2 \dots n . \gamma (\gamma + 1) \dots (\gamma + n - 1)}
\]

其中 \(\alpha, \beta, \gamma\) 是任意的实数，都不是 \(\leqslant 0\) 的整数；显然 \(u_{n}\) 从某个时刻起 \(>0\) ，或从某个时刻起 \(< 0\) 。有

\[
\begin{array}{l} \frac {u _ {n + 1}}{u _ {n}} = \frac {(\alpha + n) (\beta + n)}{(n + 1) (\gamma + n)} \\ \qquad = \left(1 + \frac {\alpha + \beta}{n} + \frac {\alpha \beta}{n ^ {2}}\right) \left(1 + \frac {\gamma + 1}{n} + \frac {\gamma}{n ^ {2}}\right) ^ {- 1} \\ \qquad = 1 + \frac {\alpha + \beta - \gamma - 1}{n} \\ \qquad + \frac {\alpha \beta - (\alpha + \beta) (\gamma + 1) + \gamma^ {2} + \gamma + 1}{n ^ {2}} + o \left(\frac {1}{n ^ {2}}\right). \end{array}
\]

拉阿贝判别法表明：级数当 \(\alpha +\beta <  \gamma\) 时收敛，当 \(\alpha +\beta >\gamma\) 时和为无穷；当 \(\alpha +\beta = \gamma\) 时，如命题 8 所示，级数的和也为无穷。

注. 1) 作为拉阿贝判别法的特例可以看出：若 \(\limsup_{n\to \infty}u_{n + 1} / u_n\) \(< 1\) ，则级数 \((u_{n})\) 收敛；反之，若 \(\liminf_{n\to \infty}u_{n + 1} / u_n > 1\) ，则级数 \((u_{n})\) 的和为无穷（达朗贝尔判别法）。

\begin{enumerate}
\def\labelenumi{\arabic{enumi})}
\setcounter{enumi}{1}
\tightlist
\item
  第二类判别法只能应用于这样的级数：其通项当 n 趋于 \(+\infty\) 时表现非常有规律；换句话说，它们的适用范围比对数判别法狭窄得多，把它们用到特别适合它们的特殊情形之外将是一个错误。例如，对级数 \((u_{n})\) （由 \(u_{2m} = 2^{-m}\) 、 \(u_{2m + 1} = 3^{-m}\) 定义），有 \(u_{2m + 1} / u_{2m} = \left(\frac{2}{3}\right)^{m}\) ， \(u_{2m + 2} / u_{2m + 1} = \frac{1}{2}\left(\frac{3}{2}\right)^{m}\) ；第一个比值趋于 0，第二个趋于 \(+\infty\) （当 \(m\) 无限增大时），因此没有任何第二类判别法可以应用；尽管如此，由于 \(u_{n} \preccurlyeq 2^{-n / 2}\) ，立即可以看出级数收敛。
\end{enumerate}

即使当 \(u_{n+1} / u_n\) 有简单的表达式时，直接估计 \(u_n\) 的主部也常常与第二类判别法一样快地得出结果。例如，对超几何级数，斯特林公式立即表明 \(u_n \sim a n^{\alpha + \beta - \gamma - 1}\) （其中 \(a\) 是一个 \(\neq 0\) 的常数），于是 0 阶对数判别法便可应用。

\appendix
\renewcommand{\thechapter}{}
\renewcommand{\thesection}{\arabic{section}}
